% !PRINTED 120
\weberpage{120}
%<<CONT

\[
 2y=u+v+w
\]
and obtain, if we write, for abbreviation,
\begin{equation}
 s=u^2+v^2+w^2,\qquad
 t=v^2w^2+w^2u^2+u^2v^2,
\tag{2}
\end{equation}
then
\[
 4y^2=s+2(vw+wu+uv),
\]
\[
 16y^4=s^2+4s(vw+wu+uv)+4t+8uvw(u+v+w).
\]
Substitution in (1) gives
\[
 s^2+4t+4as+16c+8(uvw+b)(u+v+w)
  +4(s+2a)(vw+wu+uv)=0.
\]
This equation is satisfied by the assumptions
\[
 s+2a=0,\qquad uvw+b=0,\qquad s^2+4t+4as+16c=0.
\]
With the help of the first of these three equations the third becomes
\[
 t=a^2-4c,
\]
and, putting back the values (2) for $s,t$,
\begin{equation}
\begin{array}{rcl}
 u^2+v^2+w^2&=&-2a,\\
 v^2w^2+w^2u^2+u^2v^2&=&a^2-4c,\\
 uvw&=&-b.
\end{array}
\tag{3}
\end{equation}
By \S\,7 these equations are satisfied if and only if $u^2,v^2,w^2$ are the roots of the cubic equation
\begin{equation}
 z^3+2az^2+(a^2-4c)z-b^2=0,
\tag{4}
\end{equation}
and if the signs of $u,v,w$ are so chosen that the last of equations (3) is satisfied. This last condition still allows four different determinations of signs, so that the four roots of the biquadratic equation are obtained in the following way:
\begin{equation}
\begin{array}{rcl}
2\alpha&=&u+v+w,\\
2\beta&=&u-v-w,\\
2\gamma&=&-u+v-w,\\
2\delta&=&-u-v+w.
\end{array}
\tag{5}
\end{equation}
Thus the solution of the biquadratic equation has been reduced to that of the cubic equation (4).

This equation is called a \gesperrt{cubic resolvent} of the biquadratic equation.
